\subsection{Four-window projection spectra and the exact polynomial-flag sharpness calculation}\label{endpoint-four_window_sharpness--four-window-projection-spectra-and-the-exact-polynomial-flag-sharpness-calculation}

13 September 2026. Independent written calculation for the common finite source in the arithmetic determinant transport. The original monic polynomial, coefficient frame, relation maps, and masses remain explicit.

The source \texttt{work/tau\_arithmetic\_determinant\_transport\_20260913.tex} was consulted through its first 190 lines, including AT1--AT10. Its SHA-256 at consultation was \texttt{24cf7ce5be6a43cc656641b516185df10b2e4d7824600ece657120d1c549deae}. The present calculation proves the projection statements below directly. It does not edit that source, run Lean, or assert that an arbitrary positive coefficient Gram is an arithmetic moment Gram.

\subsubsection{1. Original flag spaces and exact spectra}\label{endpoint-four_window_sharpness--original-flag-spaces-and-exact-spectra}

Fix the original monic polynomial ENDPOINTSOURCEE7E2F5FD05BDF865D519AAC0SLOT0END of degree ENDPOINTSOURCEE7E2F5FD05BDF865D519AAC0SLOT1END. Work on

ENDPOINTSOURCEE7E2F5FD05BDF865D519AAC0SLOT2END

The inner product is conjugate-linear in its first slot. Let ENDPOINTSOURCEE7E2F5FD05BDF865D519AAC0SLOT3END be the ENDPOINTSOURCEE7E2F5FD05BDF865D519AAC0SLOT4END-orthogonal projection of ENDPOINTSOURCEE7E2F5FD05BDF865D519AAC0SLOT5END onto ENDPOINTSOURCEE7E2F5FD05BDF865D519AAC0SLOT6END. Let ENDPOINTSOURCEE7E2F5FD05BDF865D519AAC0SLOT7END be the ENDPOINTSOURCEE7E2F5FD05BDF865D519AAC0SLOT8END-orthogonal projection onto the literal relation space ENDPOINTSOURCEE7E2F5FD05BDF865D519AAC0SLOT9END, with ENDPOINTSOURCEE7E2F5FD05BDF865D519AAC0SLOT10END. In particular ENDPOINTSOURCEE7E2F5FD05BDF865D519AAC0SLOT11END. Define

ENDPOINTSOURCEE7E2F5FD05BDF865D519AAC0SLOT12END

All projections in each one of the two displayed flags commute: if ENDPOINTSOURCEE7E2F5FD05BDF865D519AAC0SLOT13END, their orthogonal projections satisfy ENDPOINTSOURCEE7E2F5FD05BDF865D519AAC0SLOT14END, as follows by decomposing into ENDPOINTSOURCEE7E2F5FD05BDF865D519AAC0SLOT15END.

Write

ENDPOINTSOURCEE7E2F5FD05BDF865D519AAC0SLOT16END

Their dimensions are ENDPOINTSOURCEE7E2F5FD05BDF865D519AAC0SLOT17END, respectively, because multiplication by the nonzero polynomial ENDPOINTSOURCEE7E2F5FD05BDF865D519AAC0SLOT18END is injective. The exact eigenspaces of ENDPOINTSOURCEE7E2F5FD05BDF865D519AAC0SLOT19END are

ENDPOINTSOURCEE7E2F5FD05BDF865D519AAC0SLOT20END

Every direct sum in (FS4) is orthogonal. Applying the three projections of (FS2) to each of these spaces gives the stated eigenvalue, and the dimensions sum to ENDPOINTSOURCEE7E2F5FD05BDF865D519AAC0SLOT21END. For ENDPOINTSOURCEE7E2F5FD05BDF865D519AAC0SLOT22END, ENDPOINTSOURCEE7E2F5FD05BDF865D519AAC0SLOT23END, the last space is zero, and ENDPOINTSOURCEE7E2F5FD05BDF865D519AAC0SLOT24END.

The exact eigenspaces of ENDPOINTSOURCEE7E2F5FD05BDF865D519AAC0SLOT25END are

ENDPOINTSOURCEE7E2F5FD05BDF865D519AAC0SLOT26END

This follows by the same nested decomposition of the degree flag. Thus both operators have precisely the spectrum

ENDPOINTSOURCEE7E2F5FD05BDF865D519AAC0SLOT27END

The notation ENDPOINTSOURCEE7E2F5FD05BDF865D519AAC0SLOT28END means an eigenvalue ENDPOINTSOURCEE7E2F5FD05BDF865D519AAC0SLOT29END of multiplicity ENDPOINTSOURCEE7E2F5FD05BDF865D519AAC0SLOT30END. No claim that the spectrum of their difference is fixed by these two individual spectra is made.

\subsubsection{2. Best bound from the two individual spectra alone}\label{endpoint-four_window_sharpness--best-bound-from-the-two-individual-spectra-alone}

Let ENDPOINTSOURCEE7E2F5FD05BDF865D519AAC0SLOT31END be any ENDPOINTSOURCEE7E2F5FD05BDF865D519AAC0SLOT32END-self-adjoint endomorphism with eigenvalues

ENDPOINTSOURCEE7E2F5FD05BDF865D519AAC0SLOT33END

For a rank-ENDPOINTSOURCEE7E2F5FD05BDF865D519AAC0SLOT34END orthogonal projection ENDPOINTSOURCEE7E2F5FD05BDF865D519AAC0SLOT35END, choose an orthonormal eigenbasis ENDPOINTSOURCEE7E2F5FD05BDF865D519AAC0SLOT36END of ENDPOINTSOURCEE7E2F5FD05BDF865D519AAC0SLOT37END, and put ENDPOINTSOURCEE7E2F5FD05BDF865D519AAC0SLOT38END. Then ENDPOINTSOURCEE7E2F5FD05BDF865D519AAC0SLOT39END, ENDPOINTSOURCEE7E2F5FD05BDF865D519AAC0SLOT40END, and

ENDPOINTSOURCEE7E2F5FD05BDF865D519AAC0SLOT41END

Transferring any positive weight from a lower coefficient to an available higher coefficient does not decrease this sum. Repeating gives the exact two-sided inequality

ENDPOINTSOURCEE7E2F5FD05BDF865D519AAC0SLOT42END

If the eigenvalues of ENDPOINTSOURCEE7E2F5FD05BDF865D519AAC0SLOT43END are distinct, equality on either side forces the corresponding spectral projection: a non-extremal weight transfer would otherwise strictly improve the sum, and a diagonal projection value zero or one forces the corresponding basis vector into its kernel or image.

For any self-adjoint operator with spectrum (FS6), its spectral projections satisfy

ENDPOINTSOURCEE7E2F5FD05BDF865D519AAC0SLOT44END

Applying (FS8) twice gives

ENDPOINTSOURCEE7E2F5FD05BDF865D519AAC0SLOT45END

Subtracting, and then applying the same reasoning with ENDPOINTSOURCEE7E2F5FD05BDF865D519AAC0SLOT46END interchanged, proves

ENDPOINTSOURCEE7E2F5FD05BDF865D519AAC0SLOT47END

For ENDPOINTSOURCEE7E2F5FD05BDF865D519AAC0SLOT48END, the sum is empty and the right side is exactly ENDPOINTSOURCEE7E2F5FD05BDF865D519AAC0SLOT49END.

This is the best bound that uses only the individual spectra. To prove sharpness in that class, use the eigenbasis of ENDPOINTSOURCEE7E2F5FD05BDF865D519AAC0SLOT50END, put

ENDPOINTSOURCEE7E2F5FD05BDF865D519AAC0SLOT51END

Both have spectrum (FS6), and direct multiplication gives ENDPOINTSOURCEE7E2F5FD05BDF865D519AAC0SLOT52END. Swapping them gives the negative equality. Their difference has the ordered spectrum

ENDPOINTSOURCEE7E2F5FD05BDF865D519AAC0SLOT53END

These are exact independent-unitary-orbit equality fixtures. They do not yet assert realization by the two literal polynomial flags in (FS1)--(FS3).

\subsubsection{3. Why the literal polynomial flags cannot attain nontrivial equality}\label{endpoint-four_window_sharpness--why-the-literal-polynomial-flags-cannot-attain-nontrivial-equality}

Two exact polynomial identities govern the flag positions:

ENDPOINTSOURCEE7E2F5FD05BDF865D519AAC0SLOT54END

The first follows because a nonzero multiple of a degree-ENDPOINTSOURCEE7E2F5FD05BDF865D519AAC0SLOT55END polynomial has degree at least ENDPOINTSOURCEE7E2F5FD05BDF865D519AAC0SLOT56END. For the second, ENDPOINTSOURCEE7E2F5FD05BDF865D519AAC0SLOT57END: a nonconstant multiplier has degree above ENDPOINTSOURCEE7E2F5FD05BDF865D519AAC0SLOT58END. Thus the sum has dimension ENDPOINTSOURCEE7E2F5FD05BDF865D519AAC0SLOT59END, and is the whole space.

Let ENDPOINTSOURCEE7E2F5FD05BDF865D519AAC0SLOT60END, ENDPOINTSOURCEE7E2F5FD05BDF865D519AAC0SLOT61END, ENDPOINTSOURCEE7E2F5FD05BDF865D519AAC0SLOT62END, and ENDPOINTSOURCEE7E2F5FD05BDF865D519AAC0SLOT63END. Equations (FS4)--(FS5) and (FS14) imply

ENDPOINTSOURCEE7E2F5FD05BDF865D519AAC0SLOT64END

For the first, ENDPOINTSOURCEE7E2F5FD05BDF865D519AAC0SLOT65END whereas ENDPOINTSOURCEE7E2F5FD05BDF865D519AAC0SLOT66END. For the second, a vector in ENDPOINTSOURCEE7E2F5FD05BDF865D519AAC0SLOT67END is orthogonal to both ENDPOINTSOURCEE7E2F5FD05BDF865D519AAC0SLOT68END and ENDPOINTSOURCEE7E2F5FD05BDF865D519AAC0SLOT69END; their sum is the whole space by (FS14).

Here is the requested simple-spectrum obstruction. If equality held in the positive version of (FS11) with distinct ENDPOINTSOURCEE7E2F5FD05BDF865D519AAC0SLOT70END, both bounds in (FS10) would have to be equalities. The ENDPOINTSOURCEE7E2F5FD05BDF865D519AAC0SLOT71END eigenspace would be the span of the highest ENDPOINTSOURCEE7E2F5FD05BDF865D519AAC0SLOT72END eigenvectors of ENDPOINTSOURCEE7E2F5FD05BDF865D519AAC0SLOT73END, and the kernel of ENDPOINTSOURCEE7E2F5FD05BDF865D519AAC0SLOT74END would be the span of the highest ENDPOINTSOURCEE7E2F5FD05BDF865D519AAC0SLOT75END. For ENDPOINTSOURCEE7E2F5FD05BDF865D519AAC0SLOT76END, this contradicts the first equation in (FS15). For ENDPOINTSOURCEE7E2F5FD05BDF865D519AAC0SLOT77END, the image of ENDPOINTSOURCEE7E2F5FD05BDF865D519AAC0SLOT78END would contain the highest eigenline, which is ENDPOINTSOURCEE7E2F5FD05BDF865D519AAC0SLOT79END. That would put the constant polynomial in ENDPOINTSOURCEE7E2F5FD05BDF865D519AAC0SLOT80END, contrary to (FS14). The negative equality is treated by applying the same argument to ENDPOINTSOURCEE7E2F5FD05BDF865D519AAC0SLOT81END.

In fact the inequality is strict for every nonscalar ENDPOINTSOURCEE7E2F5FD05BDF865D519AAC0SLOT82END, even when its spectrum has repeated eigenvalues. A full proof follows, retaining the spectral gaps rather than assuming they are nonzero individually.

For ENDPOINTSOURCEE7E2F5FD05BDF865D519AAC0SLOT83END, let ENDPOINTSOURCEE7E2F5FD05BDF865D519AAC0SLOT84END be the sum of the largest ENDPOINTSOURCEE7E2F5FD05BDF865D519AAC0SLOT85END entries in (FS13). Applying (FS8) to ENDPOINTSOURCEE7E2F5FD05BDF865D519AAC0SLOT86END and ENDPOINTSOURCEE7E2F5FD05BDF865D519AAC0SLOT87END, with a rank-ENDPOINTSOURCEE7E2F5FD05BDF865D519AAC0SLOT88END projection ENDPOINTSOURCEE7E2F5FD05BDF865D519AAC0SLOT89END as the variable, gives

ENDPOINTSOURCEE7E2F5FD05BDF865D519AAC0SLOT90END

Equality is impossible in every nontrivial rank:

\begin{itemize}
\tightlist
\item
  For ENDPOINTSOURCEE7E2F5FD05BDF865D519AAC0SLOT91END, equality requires ENDPOINTSOURCEE7E2F5FD05BDF865D519AAC0SLOT92END and ENDPOINTSOURCEE7E2F5FD05BDF865D519AAC0SLOT93END, contradicting (FS15).
\item
  For ENDPOINTSOURCEE7E2F5FD05BDF865D519AAC0SLOT94END and ENDPOINTSOURCEE7E2F5FD05BDF865D519AAC0SLOT95END, minimization of ENDPOINTSOURCEE7E2F5FD05BDF865D519AAC0SLOT96END forces ENDPOINTSOURCEE7E2F5FD05BDF865D519AAC0SLOT97END; maximization of ENDPOINTSOURCEE7E2F5FD05BDF865D519AAC0SLOT98END forces that image to contain all of ENDPOINTSOURCEE7E2F5FD05BDF865D519AAC0SLOT99END. Again (FS15) excludes it. If ENDPOINTSOURCEE7E2F5FD05BDF865D519AAC0SLOT100END, the same equality instead forces ENDPOINTSOURCEE7E2F5FD05BDF865D519AAC0SLOT101END, excluded by (FS14).
\item
  For ENDPOINTSOURCEE7E2F5FD05BDF865D519AAC0SLOT102END, the maximum of ENDPOINTSOURCEE7E2F5FD05BDF865D519AAC0SLOT103END forces ENDPOINTSOURCEE7E2F5FD05BDF865D519AAC0SLOT104END, since this is its complete positive spectral space. The minimum of ENDPOINTSOURCEE7E2F5FD05BDF865D519AAC0SLOT105END forces ENDPOINTSOURCEE7E2F5FD05BDF865D519AAC0SLOT106END to contain ENDPOINTSOURCEE7E2F5FD05BDF865D519AAC0SLOT107END. Their intersection in (FS14) is zero, so this is impossible.
\item
  For ENDPOINTSOURCEE7E2F5FD05BDF865D519AAC0SLOT108END, use the complementary projection ENDPOINTSOURCEE7E2F5FD05BDF865D519AAC0SLOT109END, whose rank is between one and ENDPOINTSOURCEE7E2F5FD05BDF865D519AAC0SLOT110END. Since ENDPOINTSOURCEE7E2F5FD05BDF865D519AAC0SLOT111END, equality would require ENDPOINTSOURCEE7E2F5FD05BDF865D519AAC0SLOT112END, excluded by (FS15).
\end{itemize}

These equality conditions use only the spectral levels zero, one, and two of ENDPOINTSOURCEE7E2F5FD05BDF865D519AAC0SLOT113END; their eigenspace descriptions follow by the same projection-weight argument proving (FS8).

The set of rank-ENDPOINTSOURCEE7E2F5FD05BDF865D519AAC0SLOT114END orthogonal projections is compact in finite dimension: it is closed and bounded in the matrix space by the equations ENDPOINTSOURCEE7E2F5FD05BDF865D519AAC0SLOT115END, ENDPOINTSOURCEE7E2F5FD05BDF865D519AAC0SLOT116END, and ENDPOINTSOURCEE7E2F5FD05BDF865D519AAC0SLOT117END. Hence the maximum in (FS16) is attained, and its failure to equal ENDPOINTSOURCEE7E2F5FD05BDF865D519AAC0SLOT118END proves strict inequality for that maximum. If ENDPOINTSOURCEE7E2F5FD05BDF865D519AAC0SLOT119END are the actual eigenvalues of ENDPOINTSOURCEE7E2F5FD05BDF865D519AAC0SLOT120END, (FS8) applied to ENDPOINTSOURCEE7E2F5FD05BDF865D519AAC0SLOT121END identifies this maximum. Therefore

ENDPOINTSOURCEE7E2F5FD05BDF865D519AAC0SLOT122END

For completeness, trace rearrangement for two arbitrary self-adjoint matrices follows by the same projection argument. In an eigenbasis of ENDPOINTSOURCEE7E2F5FD05BDF865D519AAC0SLOT123END, write

ENDPOINTSOURCEE7E2F5FD05BDF865D519AAC0SLOT124END

where ENDPOINTSOURCEE7E2F5FD05BDF865D519AAC0SLOT125END projects onto its highest ENDPOINTSOURCEE7E2F5FD05BDF865D519AAC0SLOT126END eigenvectors, with any fixed choice inside a repeated eigenspace. The trace of the scalar term against ENDPOINTSOURCEE7E2F5FD05BDF865D519AAC0SLOT127END is zero, so (FS17) gives

ENDPOINTSOURCEE7E2F5FD05BDF865D519AAC0SLOT128END

whenever ENDPOINTSOURCEE7E2F5FD05BDF865D519AAC0SLOT129END is nonscalar. At least one spectral gap is then positive, so the second inequality is strict. Applying the same proof to ENDPOINTSOURCEE7E2F5FD05BDF865D519AAC0SLOT130END gives

ENDPOINTSOURCEE7E2F5FD05BDF865D519AAC0SLOT131END

For scalar ENDPOINTSOURCEE7E2F5FD05BDF865D519AAC0SLOT132END, both sides of (FS11) are zero and equality holds. This completely identifies actual finite positive-metric equality, while allowing the sharp supremum described next.

\subsubsection{4. Explicit positive metrics approaching the bound on the unchanged polynomial flags}\label{endpoint-four_window_sharpness--explicit-positive-metrics-approaching-the-bound-on-the-unchanged-polynomial-flags}

Keep the same original monic ENDPOINTSOURCEE7E2F5FD05BDF865D519AAC0SLOT133END. The following ordered list is a basis of ENDPOINTSOURCEE7E2F5FD05BDF865D519AAC0SLOT134END:

ENDPOINTSOURCEE7E2F5FD05BDF865D519AAC0SLOT135END

Its elements have strictly increasing degrees with leading coefficient one. Thus their first ENDPOINTSOURCEE7E2F5FD05BDF865D519AAC0SLOT136END spans equal the original ENDPOINTSOURCEE7E2F5FD05BDF865D519AAC0SLOT137END. Neither ENDPOINTSOURCEE7E2F5FD05BDF865D519AAC0SLOT138END nor its coefficients have been changed.

Choose any retained positive mass ENDPOINTSOURCEE7E2F5FD05BDF865D519AAC0SLOT139END and any ENDPOINTSOURCEE7E2F5FD05BDF865D519AAC0SLOT140END. Let ENDPOINTSOURCEE7E2F5FD05BDF865D519AAC0SLOT141END be the standard coordinate vectors, and define the actual linear isomorphism

ENDPOINTSOURCEE7E2F5FD05BDF865D519AAC0SLOT142END

Its diagonal in the ordered basis is ENDPOINTSOURCEE7E2F5FD05BDF865D519AAC0SLOT143END on the first ENDPOINTSOURCEE7E2F5FD05BDF865D519AAC0SLOT144END entries and ENDPOINTSOURCEE7E2F5FD05BDF865D519AAC0SLOT145END on the last ENDPOINTSOURCEE7E2F5FD05BDF865D519AAC0SLOT146END, so its determinant is ENDPOINTSOURCEE7E2F5FD05BDF865D519AAC0SLOT147END. Its inverse sends

ENDPOINTSOURCEE7E2F5FD05BDF865D519AAC0SLOT148END

Prescribe the positive coefficient metric by

ENDPOINTSOURCEE7E2F5FD05BDF865D519AAC0SLOT149END

The map ENDPOINTSOURCEE7E2F5FD05BDF865D519AAC0SLOT150END is its displayed isometry to the Euclidean space. In the original polynomial coefficients ENDPOINTSOURCEE7E2F5FD05BDF865D519AAC0SLOT151END, its complete norm is

ENDPOINTSOURCEE7E2F5FD05BDF865D519AAC0SLOT152END

In particular ENDPOINTSOURCEE7E2F5FD05BDF865D519AAC0SLOT153END, ENDPOINTSOURCEE7E2F5FD05BDF865D519AAC0SLOT154END, and ENDPOINTSOURCEE7E2F5FD05BDF865D519AAC0SLOT155END. Every unlisted cross pairing between the distinct two-coordinate blocks and ENDPOINTSOURCEE7E2F5FD05BDF865D519AAC0SLOT156END is zero. Each two-coordinate block has Gram

ENDPOINTSOURCEE7E2F5FD05BDF865D519AAC0SLOT157END

whose determinant is ENDPOINTSOURCEE7E2F5FD05BDF865D519AAC0SLOT158END. The total Gram determinant in (FS20) is exactly ENDPOINTSOURCEE7E2F5FD05BDF865D519AAC0SLOT159END. As (FS20) has determinant one relative to the increasing monomial basis, that is also the determinant there. The mass ENDPOINTSOURCEE7E2F5FD05BDF865D519AAC0SLOT160END, all cross terms, and all powers of ENDPOINTSOURCEE7E2F5FD05BDF865D519AAC0SLOT161END remain present.

Because (FS21) is triangular by degree, every transformed degree flag is the standard coordinate flag. Hence ENDPOINTSOURCEE7E2F5FD05BDF865D519AAC0SLOT162END, in these isometric coordinates, has value zero on ENDPOINTSOURCEE7E2F5FD05BDF865D519AAC0SLOT163END, value one on ENDPOINTSOURCEE7E2F5FD05BDF865D519AAC0SLOT164END, and value two on ENDPOINTSOURCEE7E2F5FD05BDF865D519AAC0SLOT165END. Put

ENDPOINTSOURCEE7E2F5FD05BDF865D519AAC0SLOT166END

These vectors are orthonormal and orthogonal to ENDPOINTSOURCEE7E2F5FD05BDF865D519AAC0SLOT167END. The transformed relation flag has ENDPOINTSOURCEE7E2F5FD05BDF865D519AAC0SLOT168END, ENDPOINTSOURCEE7E2F5FD05BDF865D519AAC0SLOT169END, and ENDPOINTSOURCEE7E2F5FD05BDF865D519AAC0SLOT170END. Thus the exact transformed operator is

ENDPOINTSOURCEE7E2F5FD05BDF865D519AAC0SLOT171END

Here the normalization in (FS26) states an orthonormal coordinate map inside the specified metric; the original norms remain (FS24)--(FS25).

\subsubsection{5. Fixed-spectrum self-adjoint control and its exact trace}\label{endpoint-four_window_sharpness--fixed-spectrum-self-adjoint-control-and-its-exact-trace}

For the prescribed ordered numbers in (FS7), define a diagonal operator ENDPOINTSOURCEE7E2F5FD05BDF865D519AAC0SLOT172END in the transformed coordinates by

ENDPOINTSOURCEE7E2F5FD05BDF865D519AAC0SLOT173END

This lists each eigenvalue exactly once. Define ENDPOINTSOURCEE7E2F5FD05BDF865D519AAC0SLOT174END, an endomorphism of the original polynomial space. Equation (FS23) proves it is ENDPOINTSOURCEE7E2F5FD05BDF865D519AAC0SLOT175END-self-adjoint with exactly those eigenvalues. Its polynomial-coordinate formula is also explicit. Set

ENDPOINTSOURCEE7E2F5FD05BDF865D519AAC0SLOT176END

Then

ENDPOINTSOURCEE7E2F5FD05BDF865D519AAC0SLOT177END

In fact this endomorphism in the original basis is independent of ENDPOINTSOURCEE7E2F5FD05BDF865D519AAC0SLOT178END. Its self-adjointness for every metric (FS25) can be verified by multiplying each block: the product of its Gram with its matrix is ENDPOINTSOURCEE7E2F5FD05BDF865D519AAC0SLOT179END, which is Hermitian.

Using (FS26)--(FS28), each of the first ENDPOINTSOURCEE7E2F5FD05BDF865D519AAC0SLOT180END two-coordinate blocks contributes

ENDPOINTSOURCEE7E2F5FD05BDF865D519AAC0SLOT181END

and the last contributes ENDPOINTSOURCEE7E2F5FD05BDF865D519AAC0SLOT182END. The ENDPOINTSOURCEE7E2F5FD05BDF865D519AAC0SLOT183END terms cancel with their equal coefficient one. Therefore the complete trace is

ENDPOINTSOURCEE7E2F5FD05BDF865D519AAC0SLOT184END

This proves that (FS11) is a sharp supremum even with the original polynomial and relation flags held fixed, when arbitrary positive coefficient metrics are allowed to vary. For nonscalar ENDPOINTSOURCEE7E2F5FD05BDF865D519AAC0SLOT185END, it is never an attained maximum at a finite positive metric by (FS19). The limit ENDPOINTSOURCEE7E2F5FD05BDF865D519AAC0SLOT186END has Gram determinant zero by (FS25), so it is not an admissible equality metric.

The case ENDPOINTSOURCEE7E2F5FD05BDF865D519AAC0SLOT187END is present literally: ENDPOINTSOURCEE7E2F5FD05BDF865D519AAC0SLOT188END, ENDPOINTSOURCEE7E2F5FD05BDF865D519AAC0SLOT189END, ENDPOINTSOURCEE7E2F5FD05BDF865D519AAC0SLOT190END, ENDPOINTSOURCEE7E2F5FD05BDF865D519AAC0SLOT191END, and (FS30) is ENDPOINTSOURCEE7E2F5FD05BDF865D519AAC0SLOT192END.

The difference operator in this family has the exact spectrum

ENDPOINTSOURCEE7E2F5FD05BDF865D519AAC0SLOT193END

To verify it, each first pair in (FS27) gives twice the difference of the rank-one projection onto ENDPOINTSOURCEE7E2F5FD05BDF865D519AAC0SLOT194END and that onto ENDPOINTSOURCEE7E2F5FD05BDF865D519AAC0SLOT195END. Its matrix has trace zero and determinant ENDPOINTSOURCEE7E2F5FD05BDF865D519AAC0SLOT196END, so its eigenvalues are ENDPOINTSOURCEE7E2F5FD05BDF865D519AAC0SLOT197END, before multiplication by two. The last pair has the same matrix without the factor two. The common vector ENDPOINTSOURCEE7E2F5FD05BDF865D519AAC0SLOT198END is in the zero eigenspace. This also shows directly how the extremal vector (FS13) is approached.

If it is desired that the chosen ENDPOINTSOURCEE7E2F5FD05BDF865D519AAC0SLOT199END be an actual logarithmic derivative of a positive metric path, retain the same coefficient frame and define

ENDPOINTSOURCEE7E2F5FD05BDF865D519AAC0SLOT200END

The first expression is positive for every real ENDPOINTSOURCEE7E2F5FD05BDF865D519AAC0SLOT201END. The equality follows from ENDPOINTSOURCEE7E2F5FD05BDF865D519AAC0SLOT202END, term by term in the exponential series. Differentiation gives ENDPOINTSOURCEE7E2F5FD05BDF865D519AAC0SLOT203END, with no rescaling of its eigenvalues. A locally affine path ENDPOINTSOURCEE7E2F5FD05BDF865D519AAC0SLOT204END gives the same derivative at zero and is positive whenever all ENDPOINTSOURCEE7E2F5FD05BDF865D519AAC0SLOT205END are positive.

\subsubsection{6. The moment-Gram class remains a stricter class}\label{endpoint-four_window_sharpness--the-moment-gram-class-remains-a-stricter-class}

The original arithmetic and Gamma norms have the additional form

ENDPOINTSOURCEE7E2F5FD05BDF865D519AAC0SLOT206END

For all polynomials for which the displayed products belong to the finite space, the literal vertical coordinate forces

ENDPOINTSOURCEE7E2F5FD05BDF865D519AAC0SLOT207END

Indeed the integrands on the left add the factor ENDPOINTSOURCEE7E2F5FD05BDF865D519AAC0SLOT208END. The original mass is the value at ENDPOINTSOURCEE7E2F5FD05BDF865D519AAC0SLOT209END. These relations are extra linear constraints on the coefficient Gram; positive definiteness by itself does not supply them.

The family (FS24) does not satisfy these moment constraints. Set ENDPOINTSOURCEE7E2F5FD05BDF865D519AAC0SLOT210END, ENDPOINTSOURCEE7E2F5FD05BDF865D519AAC0SLOT211END; then ENDPOINTSOURCEE7E2F5FD05BDF865D519AAC0SLOT212END has degree ENDPOINTSOURCEE7E2F5FD05BDF865D519AAC0SLOT213END. The exact family gives

ENDPOINTSOURCEE7E2F5FD05BDF865D519AAC0SLOT214END

If ENDPOINTSOURCEE7E2F5FD05BDF865D519AAC0SLOT215END, ENDPOINTSOURCEE7E2F5FD05BDF865D519AAC0SLOT216END and its image under (FS21) is a coordinate vector orthogonal to ENDPOINTSOURCEE7E2F5FD05BDF865D519AAC0SLOT217END; hence ENDPOINTSOURCEE7E2F5FD05BDF865D519AAC0SLOT218END. If ENDPOINTSOURCEE7E2F5FD05BDF865D519AAC0SLOT219END, monicity writes ENDPOINTSOURCEE7E2F5FD05BDF865D519AAC0SLOT220END, and ENDPOINTSOURCEE7E2F5FD05BDF865D519AAC0SLOT221END. Thus the left side of (FS34) is ENDPOINTSOURCEE7E2F5FD05BDF865D519AAC0SLOT222END for ENDPOINTSOURCEE7E2F5FD05BDF865D519AAC0SLOT223END, and ENDPOINTSOURCEE7E2F5FD05BDF865D519AAC0SLOT224END for ENDPOINTSOURCEE7E2F5FD05BDF865D519AAC0SLOT225END, whereas its right side is zero. Since ENDPOINTSOURCEE7E2F5FD05BDF865D519AAC0SLOT226END, neither case can be a moment Gram of the form (FS33), for any real vertical centre.

Consequently the family proves sharpness for arbitrary finite positive coefficient metrics and for the literal polynomial flags. It does not prove sharpness within the original arithmetic or Gamma moment-Gram class. The exact map (FS23) and the explicit violated relation (FS34) identify the relation between those classes, while preserving the original polynomial and its full relation space. No estimate on the original source multiplier, and no arithmetic equality fixture, is inferred from the larger class.
